\section{Follow one bounded task in an installed client}\label{sec:ch2-native-observation}

The files, tools and rules become easier to understand when we watch them work together. We will return to the familiar Chapter~1 defect in a fresh copy: reject malformed, nonfinite and nonpositive amounts while preserving the public API, \texttt{BR-001/v1} and the EUR\,5,000 boundary. Using a known small problem lets us concentrate on the client rather than learning another domain at the same time. Preserve your cumulative working copy; this optional observation has its own disposable copy and return point.

The primary route uses the installed \textbf{OpenCode 1.19.3} client with \textbf{Qwen3-4B Q4\_K\_M} served locally by \textbf{llama.cpp b11799}. OpenCode sends requests to the model, offers tools and applies its application permission rules. The model proposes actions; the client executes an allowed action and returns its result to the model. The reader's recorder preserves the surrounding command and file changes. It does not supply the agent's reasoning or substitute for the native client. OpenCode documents command-line execution, session export and configurable model providers; the companion pins and observes the particular release used here \parencite{opencodeCliD6_2026,opencodeProvidersD6_2026}.

Begin with \path{practice/native_client/README.md}. It gives the fixed download hashes, installation, server startup and preparation commands for the observed Linux x86-64 environment. It also lets you inspect the retained observation without installing the approximately 2.5\,GB model. The model requests in that observation used a loopback endpoint. Installation required public registries, and the client made failed background package requests during the task, so the whole host is not described as offline.

\subsection*{Step 1: establish what the client is being asked to repair}

Open the intent before starting the client. For this task, a zero amount is invalid even though it is below EUR\,5,000. A malformed decimal should produce the specified rejection rather than an exception. These facts come from the approved teaching intent, not from whatever change the model finds easiest to generate. The initial acceptance run gives the model a concrete failing observation and gives the reviewer a baseline.

The companion setup records the unchanged companion location in \texttt{ESA\_COMPANION} and defines three separate working locations: \texttt{ESA\_RUNTIME} holds the installed client and model, \texttt{ESA\_WORK} is a new Chapter~1 study copy, and \texttt{ESA\_PACKET} retains evidence outside that copy. After setup, run the following from the unchanged companion root:

\begin{lstlisting}[language=bash]
python3 -B reader_start.py new 1 \
  --destination "$ESA_WORK"
python3 -B validation/observe_native.py prepare \
  --source "$ESA_WORK" --output "$ESA_PACKET" \
  --actions R06
cd "$ESA_WORK"
python3 -B reader_side/tools/scope_guard.py --chapter CP01 \
  --capture reader_side/evidence/CP01_NATIVE_BASELINE_01.json
python3 -B reader_side/tools/run_acceptance.py \
  --output reader_side/evidence/CH01_NATIVE_INITIAL_01.json
\end{lstlisting}

Expect three passes, two failures and two errors from the prepared starter. The negative and zero cases return the wrong route; the malformed and nonfinite cases raise exceptions. The command's nonzero exit is expected here. An unrelated import failure or a changed protected file is a different problem and should be resolved before delegation. The preparation step retains the starting file bytes, while the scope baseline identifies the router as the sole source that may change.

\begin{conceptbox}{A baseline explains what a later result means}
A baseline records the state before an intervention. Here it answers two practical questions: which defect did the client receive, and which files were permitted to change? A later passing test report without that starting state cannot show whether the client repaired the intended problem, began with a worked answer or changed the expectations. A file hash identifies bytes; it does not decide whether those bytes express the right requirement.
\end{conceptbox}

\subsection*{Step 2: inspect effective permissions before making a model request}

The companion's \path{practice/native_client/configure_chapter1.py} prepares a configuration and a task prompt for the actual location of your copy. It checks the starter and baseline, detects the Git worktree used for permission paths, and refuses to overwrite an earlier preparation. It does not invoke the model. The prompt supplies exact absolute paths to the intent, tests, router, initial result and repository guidance. It asks for the bounded repair, one named acceptance command and a stop for human review.

\begin{lstlisting}[language=bash]
cd "$ESA_COMPANION"
python3 -B practice/native_client/configure_chapter1.py \
  --work "$ESA_WORK" --runtime-root "$ESA_RUNTIME" \
  --output "$ESA_PACKET/configuration"
\end{lstlisting}

Read the generated configuration before using it. Only those named study files may be read through the native read tool. Only the router may be edited. The command allowance names the exact acceptance invocation and its new output file; other commands are denied. The recorder's task declaration separately names the source and report that may change. That declaration supports a later comparison; it is not an OS permission setting.

The retained setup exposed a useful versioning failure. The current documentation used \texttt{shell}, while installed OpenCode 1.19.3 used \texttt{bash} for this tool and permission category. The earlier configuration parsed, but the effective agent reported the command tool disabled. Correcting the category to \texttt{bash} enabled the intended tool under the same narrow command rule. The reader recipe therefore captures the installed version, effective configuration and effective agent before it offers a task to the model. A successful parse is only the first check.

The recovered settings also disable automatic context compaction for this small task. \emph{Compaction} condenses accumulated conversation so that a session can continue within its usable context; it can itself make model requests. An earlier continued session began this activity after its final-looking message. The targeted configuration change makes the chosen task boundary clearer, while the process receipt still establishes when execution actually stops.

\begin{readerbox}{Read the rule as well as the tool name}
The observed effective agent enabled \texttt{read}, \texttt{edit}, \texttt{write} and \texttt{bash}. Those booleans say that the tools are offered; the ordered permission rules determine their allowed targets. An internal OpenCode tool-output directory exception also remained in the effective rules. Inspect both parts. Application allowlists and a process's operating-system rights answer different questions; this configuration does not demonstrate general filesystem or network isolation.
\end{readerbox}

\subsection*{Step 3: use the observed failure to choose the next action}

After inspecting the preparation, invoke the actual client through the generated command specification. The model, context, output and step limits are declared in that specification and its configuration. They bound the exercise; they do not guarantee that the model will repair the function.

\begin{lstlisting}[language=bash]
python3 -B validation/observe_native.py capture \
  --spec "$ESA_PACKET/configuration/CLIENT_TASK.json" \
  --output "$ESA_PACKET/attempt_01"
\end{lstlisting}

Read the native tool events before deciding what to do next. A failed file read, an unfinished model request and a rejected candidate need different corrections. Changing the code cannot fix a path that the client was never permitted to read. Granting more file access cannot make an invalid decimal comparison correct. The retained episode makes these distinctions concrete.

\begin{table}[H]\centering\small
\caption{Use an observed failure to choose a bounded recovery}\label{tab:ch2-native-recovery}
\begin{tabularx}{\textwidth}{@{}p{0.26\textwidth}XX@{}}\toprule
\textbf{Observed result} & \textbf{What it establishes} & \textbf{Decision for the next step}\\\midrule
Five reads refused & Requested paths omitted the study directory; no edit occurred & Supply the exact allowed paths and preserve the permission scope\\
Request deadline after reads & File access succeeded; the next model request stopped during prompt processing & Retain server evidence and declare an appropriate bounded CPU request budget\\
Six passes, one error & The candidate still failed the nonfinite case, despite readiness wording & Use the actual report to request correction; do not accept the summary\\
Repeated faulty insertion & Continued editing left the same error and duplicated old code & Preserve the candidates; use a fresh session and disclose more specific guidance\\
Compaction after final prose & The host continued model work after the apparent conclusion & Observe process completion and retain the additional activity\\\bottomrule
\end{tabularx}
\end{table}

The table preserves the earlier d6 failures, whose native records remain in the companion's \path{practice/native_client/observed/ch01_01/historical_d6/} extract. The code failure deserves a closer look. Parsing \texttt{NaN} as a \texttt{Decimal} can succeed, yet a later ordering comparison can raise an exception \parencite{pythonDecimalCh8}. ``It parsed'' is therefore weaker than the intent's requirement that the amount be finite and positive. The failed acceptance case directs attention to that distinction. Inspecting the diff supplies another kind of information: an unreachable copy of the old body is confusing even when a selected test does not execute it. These observations justify a focused correction without changing the approved threshold or rewriting the tests.

\begin{examplebox}{Three different results after a client stops}
\textbf{Process completion} is established by the exit status and any timeout. \textbf{Task completion} requires the actual reads, edits, checks and candidate. \textbf{Acceptance} requires the independent review and decision. The historical d6 first attempt completed the process with exit zero while leaving the task unfinished. Later, a final-looking failure message was followed by automatic context compaction in the host. Neither a reassuring exit code nor a concluding sentence supplies the other observations.
\end{examplebox}

Keep each prompt, configuration and failed result under its own attempt name. When an intervention supplies API-level guidance or starts a fresh session, record that assistance: the result then reflects those more specific instructions. A failed request without a completed usage receipt has unknown token use, and an application field containing \texttt{cost: 0} does not measure total financial cost. Chapter~14 will make this accounting distinction explicit.

\subsection*{Step 4: retain the native record and review the candidate}

The command recorder preserves stdout, stderr, literal arguments and before/after file inventories. The native session export supplies the client's own message and tool context. Use the companion's direct-file export recipe, then parse the JSON and verify the session identifier. A command returning zero does not establish that its complete output was retained. The guide preserves the collection failure that motivated this check.

After the agent stops, the reader reruns the seven acceptance cases into a new report, checks the original edit baseline and compares the candidate with the preserved starter. A green result means that those cases passed against the identified candidate. The diff still needs review against the intent: did the validation preserve the threshold, return type and policy version? The companion guide gives the commands and retains the initial sources, native events, unsuccessful attempts and independent results together under \path{practice/native_client/observed/ch01_01/}.

The fresh d7 observation on 22 October 2026 reached a useful checked candidate. Its first attempt read the router, ran acceptance before editing and retained the original three passes, two failures and two errors. An operator stopped its next model request after 555.556 seconds; no source changed. A separate guided recovery supplied the approved intent and explicit \texttt{Decimal.is\_finite()} advice before comparisons. Although its prompt requested \texttt{write}, the native client used the allowed \texttt{edit} tool to replace the function body. The 360-second outer bound ended the client after 360.108 seconds, before its planned checker or final summary completed. Independent assistant review then passed the original seven cases, fourteen additional cases planned before the candidate appeared, and the scope guard. The two case sets overlap. Their reports and the candidate hash remain separate from the native process outcome; human acceptance is still pending.

The native read records expose another context surface: reading the 671-byte router returned 9,077 characters because OpenCode appended the repository's integrated Chapters 1--14 \texttt{AGENTS.md}. Its tool metadata names the loaded file. This is direct evidence of context inclusion, not proof that the guidance caused the earlier error or timeout. The packet retains the request timings, a bounded server-log prefix and resource samples from the shared 8\,GiB host. Four completed native steps have usage receipts; two interrupted requests do not, so complete token use remains unknown. These observations support diagnosis and an honest stopping decision, not a hardware minimum or performance ranking.

This observation is not blind. The author had already seen worked solutions; the native task supplied allowed source and tests, but did not authorize reading the supplied answers or inject a repair patch. The recovery's API guidance is additional disclosed assistance. Record the same prior-exposure fact in your own exercise. Assistant review and local checks do not become human acceptance merely because they agree. Keep the human decision pending until it actually occurs. The earlier active d6 attempt 06 was excluded from its saved checkpoint; the fresh d7 result does not establish that missing outcome.

\subsection*{Step 5: connect the policy lookup with a smaller tool set}

The next task needs information rather than a code change. The client should retrieve \texttt{BR-001/v2}, request the unknown identifier \texttt{BR-999/v1}, and attempt one deliberately wrong argument type. We use the existing SDK server and catalogue from the earlier MCP example. Reading v2 makes its contents visible; it does not authorize replacing the Chapter~1 rule with that later policy.

The companion's \path{practice/native_client/configure_mcp.py} prepares a separate study and evidence directory, the actual host configuration and the exact prompt. It uses the pinned Python environment and the existing wire recorder. Only the policy lookup is offered to the model; no file-editing or shell tool is needed for this task. The preparation creates files but does not invoke a model or start the server. Inspect it before running the generated client task. A fresh preparation is needed for each launch because the wire recorder preserves its first trace rather than overwriting it.

Follow three different boundaries in the resulting record. First, the host must start the process and discover its tool. Second, a native model action must select the tool and supply arguments. Third, the server must return the actual catalogue result or refusal. A well-formed unknown identifier is a domain outcome: the tool understood the request but found no such policy. A list where the schema requires a string is an argument error. If the host rejects that call before sending an MCP request, the evidence must name that layer rather than attributing the rejection to the server.

The retained 22 October native run makes these distinctions concrete. OpenCode~1.19.3 started the canonical server with MCP SDK~2.3.0, negotiated protocol \texttt{2025-11-25}, discovered \texttt{lookup\_policy} and selected it in the model's native tool events. \texttt{BR-001/v2} returned the catalogue record; \texttt{BR-999/v1} returned \texttt{UNKNOWN\_POLICY}. The attempted wrong-type call is a useful failed claim: although the prompt requested an empty array, the native event and wire both contain the string \texttt{"["}. The server therefore returned another domain refusal. The model's first final response claimed that schema validation rejected an array. We preserve that response and reject its explanation against the actual request and result. The available evidence does not isolate whether the intended array changed in the model, provider adapter or client parsing.

A separately labelled SDK control, using the same negotiated protocol and unchanged server, really sent \texttt{[]} and received the server's string-type validation error. It is evidence about server argument handling, not an additional model-selected OpenCode call. Its first wrapper failed before starting the server because its trace directory was absent; that failure and the corrected wrapper's result remain separate. The packet thus contains actual native discovery and policy use, domain refusals, a preserved false interpretation and an independently attributed argument control. It does not claim that every diagnostic case occurred through the model.

\begin{conceptbox}{A read-only tool can run in a process with wider rights}
The lookup accepts fixed policy identifiers and exposes no file-reading or writing operation. Its Python process nevertheless runs with the rights granted by its launch environment. The optional observation declares one harmless canary file outside both the study and server directories, then instruments the serving interpreter to try reading that file at startup. This is a \emph{canary}: deliberately placed test data with no secret value, used to observe a specific boundary.

A successful read establishes access to that path by that process at that time. It does not show that the model can request arbitrary files through the policy tool. A missing probe record leaves the read unobserved. The probe changes the runtime through an explicit startup hook while leaving the canonical server source unchanged; it does not test network access, secret inheritance, external writes or isolation established later.
\end{conceptbox}

The startup probe also produced a positive result in this run: the canonical serving interpreter read the declared 50-byte harmless canary outside the study directory before initialization. The retained probe identifies that process, path and content hash. It establishes that particular startup read. It supplies no evidence of post-initialization containment, network restrictions or a model-accessible arbitrary-file operation.

Inspect the native events beside the raw MCP messages, and inspect the process probe separately. A tool label, a successful lookup and an operating-system restriction are three different observations. The same care applies to protocol versions: record the actual opening exchange rather than deriving it from an installed package's version number. The earlier SDK-only trace remains evidence about that earlier connection. The complete native collection guide retains each connection's configuration and its own resulting messages under \path{practice/native_client/observed/mcp_host_01/}.

\subsection*{A meaningful contrast without inventing a comparison result}

VS Code with GitHub Copilot remains the unobserved contrasting route. Its documented IDE workflow streams edits and normally asks the user to approve commands, while editor or administrator settings may permit automatic execution \parencite{githubCopilotAgent2026}. Our observed route runs a literal command noninteractively with one predeclared command allowance. The difference changes where the engineer inspects execution and collects evidence. It supplies no speed, quality or cost ranking. To make a paired comparison, observe both clients on separate fresh copies with the same intent, scope and review criteria.

\begin{readerbox}{A consistent packet still needs a reviewer}
The \texttt{bind} and \texttt{check} commands in \path{validation/observe_native.py} connect required evidence files to hashes and detect missing or changed records. They do not authenticate the claimed producer, certify process containment or grant acceptance. Use the native events to assess what the client did, independent checks to assess the candidate, and a named decision to record what was accepted.
\end{readerbox}

